% Generated draft: requires mathematical and source review.
\documentclass[11pt]{article}
\usepackage[margin=25mm]{geometry}
\usepackage{fontspec}
\setmainfont{Latin Modern Roman}
\usepackage{amsmath,amssymb,mathtools,unicode-math}
\setmathfont{Latin Modern Math}
\usepackage{enumitem,needspace,xurl,hyperref}
\setlength{\parindent}{0pt}
\setlength{\parskip}{5pt}
\setlength{\emergencystretch}{3em}
\newcommand{\N}{\mathbb{N}}
\newcommand{\Z}{\mathbb{Z}}
\newcommand{\Q}{\mathbb{Q}}
\newcommand{\R}{\mathbb{R}}
\newcommand{\C}{\mathbb{C}}
\newcommand{\F}{\mathbb{F}}
\newcommand{\A}{\mathbb{A}}
\newcommand{\PP}{\mathbb{P}}
\newcommand{\E}{\mathbb{E}}
\newcommand{\eps}{\varepsilon}
\newcommand{\dd}{\,\mathrm{d}}
\begin{document}
\section*{UnsolvedMath ID: 30002595}
\textbf{Source problem code:} catalogue entry
\subsection{Definitions and mathematical statement}
A knot is an embedded circle in \(S^3\), considered up to ambient isotopy. A two-bridge presentation decomposes \((S^3,K)\) along a sphere into two balls, each containing two boundary-parallel arcs. The supplied record does not specify whether the unknot is included among two-bridge knots.

For an integer \(n>1\), let \(\Sigma_n(K)\) be the connected \(n\)-fold cyclic cover of \(S^3\) branched over \(K\). Its restriction over the knot complement corresponds to the kernel of the homomorphism
\[
\pi_1(S^3\setminus K)\longrightarrow \Z/n\Z
\]
that sends a meridian to \(1\). A group \(G\) is left-orderable if it admits a strict total order \(<\) satisfying
\[
\forall g,h,a\in G,\qquad g<h\ \Longrightarrow\ ag<ah.
\]
Under this definition the trivial group is left-orderable; the supplied record gives no separate convention for that boundary case.

\textbf{Classification problem.} For every two-bridge knot \(K\) and every integer \(n>1\), determine each of the following:
\begin{enumerate}
\item whether \(\Sigma_n(K)\) is a total L-space;
\item whether \(\pi_1(\Sigma_n(K))\) is left-orderable;
\item whether \(\Sigma_n(K)\) is excellent.
\end{enumerate}
The requested classification must cover every such pair \((K,n)\). These are three requested determinations; the supplied record does not assert that they form mutually exclusive or exhaustive alternatives.

\textbf{Unresolved terminology.} The supplied record does not define ``total L-space'' or ``excellent.'' These terms are retained as source predicates, without substituting definitions or asserting equivalences with left-orderability. Their intended definitions and any associated hypotheses must be recovered from the cited source before this becomes a fully specified mathematical statement. Any rational-homology-sphere, irreducibility, orientation, or other restrictions implicit in those terms remain unresolved.
\subsection{Short English statement}
Classify the cyclic branched covers \(\Sigma_n(K)\) for every two-bridge knot \(K\) and integer \(n>1\), determining total L-space status, left-orderability of the fundamental group, and excellence. The source does not define total L-space or excellence, and its treatment of the unknot is unspecified. The supplied record labels the problem open; that status is not independently verified here.
\subsection{Sources}
\begin{itemize}
\item Local UnsolvedMath registry, supplied complete source record, permanent ID \(30002595\), title ``Classification of Cyclic Covers of Two-Bridge Knots,'' source code \texttt{catalogue entry}, in the Oberwolfach Reports Open Problems collection. No individual proposer is named; the supplied proposed year is 2014.
\item Supplied source citation: \emph{Algebraic Structures in Low-Dimensional Topology (2015)}, Oberwolfach Reports, pp.\ 1403--1458, DOI \texttt{10.4171/owr/2014/26}. \url{https://doi.org/10.4171/owr/2014/26}. The background describes these reports as documenting mathematical workshops and their open-problem discussions.
\item Supplied literature reference: ``Oberwolfach Report 26, cyclic branched covers.'' \url{https://doi.org/10.4171/owr/2014/26}.
\item Supplied literature reference: ``Hu, left-orderability and cyclic branched coverings.'' \url{https://doi.org/10.2140/agt.2015.15.399}.
\item Supplied literature reference: ``Issa--Turner, cyclic-branched-cover L-spaces.'' \url{https://doi.org/10.1112/blms.12950}.
\item Supplied literature assessment, dated exactly \texttt{2026-08-21}: the status is recorded as open; results are said to settle important subclasses, while a complete classification for all two-bridge knots and all covering degrees remains open. This is supplied dated curation derived from the OWR record, not an independently verified literature assessment. The supplied background explicitly requests independent verification of time-sensitive status and cited claims.
\item The registry labels the statement \texttt{corrected\_verified} and says that notation, typography, and wording were normalized without changing mathematical content. This records the supplied verification claim, not mathematical verification of this draft.
\end{itemize}
\end{document}
